% 学习指导第5版自测题10（光的衍射与偏振）10.1.1—10.1.10；
% 原书印刷206—207，PDF物理214—215。每题按源选项和原图检查。
\subsection{自测题10：单缝、光栅与偏振选择题}

\sourcetest{10.1.1}{206}{214}
\problemtext 单缝光垂直入射，某$\theta=30^\circ$方向
按半波带法可以把波面分成3个半波带，
求单缝宽度$a$。
\workings
单缝两边光程差$\delta=a\sin\theta$。
每个半波带对应$\lambda/2$，故
$a\sin30^\circ=3\lambda/2$，即
\[
\boxed{a=3\lambda}.
\]
选\textbf{D}。
\crosscheck 3个半波带不是3个全波长，不能漏掉$1/2$。

\sourcetest{10.1.2}{206}{214}
\problemtext 夫琅禾费单缝衍射，狭缝宽增加为原来的$3/2$，
光波波长减至原来的$3/4$，问中央明纹近轴宽度比例。
\workings
中央明纹角宽或固定焦距的焦面线宽均正比$\lambda/a$，
\[
\boxed{\frac{W'}W=
\frac{(3/4)\lambda}{(3/2)a}\bigg/
\frac{\lambda}{a}=\frac12}.
\]
选\textbf{D}。
\crosscheck 中央明纹变窄，不是随缝宽增加而变宽。

\sourcetest{10.1.3}{206}{214}
\problemtext 原书图10.1.2：单缝宽$a$略增大，
并将整个单缝沿$+y$方向微小平移；
透镜和接收屏不动，问中央衍射纹宽度及位置变化。
\workings
单缝变宽使中央衍射明纹角宽
$\Delta\theta\simeq2\lambda/a$减小。
理想平面波与透镜后焦面构成傅里叶光学布置，
平移孔径仅使衍射场的复振幅多一个线性相位因子，
其强度模不变，中央亮纹不因单缝整体平移而移动。
所以\textbf{变窄，不移动}，选\textbf{C}。
\crosscheck 必须区分源题的“把缝宽稍微\emph{变宽}”
与“缝沿$+y$小位移”；不能把两种变化混作宽度减小。

\sourcetest{10.1.4}{206}{214}
\problemtext 已有光栅在屏幕上只能显示零级和正负一级主极大，
若想出现更高级主极大，应换大还是小的光栅常数？
\workings
光栅条件$d\sin\theta=m\lambda$，
传播方向满足$|m|\le d/\lambda$。
要增大允许最高级数$m_{\max}$，
需要更大的光栅常量$d$，
故选\textbf{B}。
\crosscheck 更大的刻线密度意味着更小的$d$，
反而减少可传播级次；“高分辨刻线多”与“高阶可见”是不同优化目标。

\sourcetest{10.1.5}{206}{214}
\problemtext 多缝光栅的所有偶数级主极大均遇单缝零点而缺失。
透光缝宽$a$与相邻不透明缝间隔$b$有何关系？
\workings
主极大$d\sin\theta=m\lambda$，
单缝暗纹$a\sin\theta=k\lambda$，
缺级当$m=(d/a)k$。
所有偶级缺失而奇级保留对应$d/a=2$，
即$d=a+b=2a$，所以
\[
\boxed{a=b}.
\]
选\textbf{A}。
\crosscheck “所有偶级”指非零偶数级；
零级中心主极大不在单缝零点。

\sourcetest{10.1.6}{207}{215}
\problemtext 强度$I_0$自然光依次通过两片理想偏振片，
透振轴夹角30°，求透射强度。
\workings
通过第一片：$I_1=I_0/2$。
再按马吕斯定律：
\[
\boxed{I=I_1\cos^230^\circ
=\frac{I_0}{2}\cdot\frac34=\frac{3I_0}{8}}.
\]
选\textbf{E}。
\crosscheck 首片透射因子$1/2$必不可少。

\sourcetest{10.1.7}{207}{215}
\problemtext 三片偏振片$P_1,P_2,P_3$，$P_1$与$P_3$互相正交。
自然光照射后出射光强为$I_0/8$。
只转动中间片$P_2$使强度归零，问最小转角。
\workings
设中片轴与$P_1$轴夹角$\theta$，$P_3$轴与$P_1$夹90°，
\[
 I=\frac{I_0}{2}\cos^2\theta\sin^2\theta
 =\frac{I_0}{8}\sin^22\theta.
\]
已知$I=I_0/8$，故$\sin^22\theta=1$，
中片轴在45°模90°位置。
要归零必须使其与某一片轴平行，最小转动
\[
\boxed{\Delta\theta_{\min}=45^\circ},
\]
选\textbf{B}。
\crosscheck 纯自然光通过首片后才能应用第二、第三片的马吕斯定律。

\sourcetest{10.1.8}{207}{215}
\problemtext 入射束混合自然光$I_n$与线偏振光$I_p$，
用理想旋转分析偏振片测得最大透射强度为最小的5倍。
求$I_n/I_p$。
\workings
自然光无论分析器转角，总有$I_n/2$通过；
线偏振分量平行分析器时全部通过、垂直时不通过。
于是
\[
\frac{I_{\max}}{I_{\min}}
=\frac{I_n/2+I_p}{I_n/2}=5
\quad\Longrightarrow\quad I_p=2I_n.
\]
故
\[
\boxed{\frac{I_n}{I_p}=\frac12}.
\]
选\textbf{A}。
\crosscheck 不应把自然光本身的半强度因子错施加到已完全线偏振的分量。

\sourcetest{10.1.9}{207}{215}
\problemtext 自然光以60°从空气入射透明介质，
测得反射光完全线偏振，求折射角以及折射光偏振性质。
\workings
布儒斯特条件要求
$i_B+r_B=90^\circ$，
所以
\[
\boxed{r_B=90^\circ-60^\circ=30^\circ}.
\]
入射自然光中平行入射面的$p$分量此时不反射，
但会折射入介质；$s$分量亦部分透射。
因此折射光一般仍为\textbf{部分偏振光}。
选\textbf{B}。
\crosscheck 完全偏振的是\emph{反射}光，而非折射光。

\sourcetest{10.1.10}{207}{215}
\problemtext 自然光以布儒斯特角从空气入射玻璃，
判断反射光电矢量的振动方向。
\workings
菲涅耳反射系数$p$分量在布儒斯特角为零，
$s$分量保留，故反射光\textbf{完全线偏振，电场垂直入射面}。
选\textbf{C}。
\crosscheck 题目问“振动方向”是指光电场方向，
不是光线传播方向；不能把两者混淆。
